\documentclass[memoire, 12pt]{report}
\usepackage[top = 1.9cm, bottom = 1.5cm, left = 1.9cm, right = 2.1cm]{geometry}
\usepackage{graphicx} 
\usepackage{enumitem}
\usepackage{multicol}
\usepackage{tabto}
\usepackage{multirow}
\usepackage{multibib}
\usepackage{multirow}
\usepackage{tabularx}
\newcites{biblio}{Bibliographie}
\newcites{other}{Autres r\'ef\'erences}
\usepackage{amssymb}
\usepackage{amsmath}
\usepackage{graphicx}
\usepackage{amsfonts}
\usepackage{lmodern}
\usepackage{caption}
\usepackage{subcaption}
\usepackage[babel=true]{csquotes}
\setlength{\fboxrule}{0.01cm}
\setlength{\fboxsep}{0.5cm}
\usepackage{array}
\usepackage{tikz}
\usepackage{lipsum}
\usepackage{setspace}
\usepackage{ragged2e}
\usepackage{url}
\usepackage{float}
\usepackage{pdfpages}
\usepackage{rotating}
\usepackage{glossaries}
%\usepackage[thinlines]{easytable}
\usepackage{hyperref}
\usepackage[export]{adjustbox}
\usepackage[bottom]{footmisc}
%\usepackage{algpseudocode}
\usepackage{algorithm}
\usepackage{algorithmic}
\usepackage[normalem]{ulem}
\useunder{\uline}{\ul}{}
\usepackage{glossaries}
\usepackage{listings}
\usepackage{xcolor}
\usepackage{minted}

\usepackage{array}
\usepackage{longtable}
\usepackage[table,xcdraw]{xcolor}

\usepackage[utf8]{inputenc}   % encodage du fichier source
\usepackage[T1]{fontenc}      % encodage des polices
\usepackage[french]{babel}    % pour le français

% Configuration des styles pour le code Python

\definecolor{codegreen}{rgb}{0,0.6,0}
\definecolor{codegray}{rgb}{0.5,0.5,0.5}
\definecolor{codepurple}{rgb}{0.58,0,0.82}
\definecolor{backcolour}{rgb}{0.95,0.95,0.92}

\lstdefinestyle{python}{
    backgroundcolor=\color{backcolour},   
    commentstyle=\color{codegreen},
    keywordstyle=\color{magenta},
    numberstyle=\tiny\color{codegray},
    stringstyle=\color{codepurple},
    basicstyle=\ttfamily\footnotesize,
    breakatwhitespace=false,         
    breaklines=true,                 
    captionpos=b,                    
    keepspaces=true,                 
    numbers=left,                    
    numbersep=5pt,                  
    showspaces=false,                
    showstringspaces=false,
    showtabs=false,                  
    tabsize=2
}

\lstset{style=python}



\usepackage[utf8]{inputenc}  
\usepackage[T1]{fontenc} 
%\usepackage{fancyhdr}
\usepackage[Conny]{fncychap}
%Conny
%Bjornstrup
%\pagestyle{Conny}
\usepackage[french]{babel}
%\renewcommand{\footrulewidth}{3pt}
\makeglossaries
\title{TAF_N1_NGWAMBE MARIELLA}
\author{}
\date{MOIS_ICI 2025}

\begin{document}
\begin{titlepage}

	\begin{tikzpicture}[remember picture,overlay,inner sep=0,outer sep=0]
		\draw[orange!90!orange,line width=4pt] ([xshift=-1.5cm,yshift=-2cm]current page.north east) coordinate (A)--([xshift=1.5cm,yshift=-2cm]current page.north west) coordinate(B)--([xshift=1.5cm,yshift=2cm]current page.south west) coordinate (C)--([xshift=-1.5cm,yshift=2cm]current page.south east) coordinate(D)--cycle;
		
		\draw ([yshift=0.5cm,xshift=-0.5cm]A)-- ([yshift=0.5cm,xshift=0.5cm]B)--
		([yshift=-0.5cm,xshift=0.5cm]B) --([yshift=-0.5cm,xshift=-0.5cm]B)--([yshift=0.5cm,xshift=-0.5cm]C)--([yshift=0.5cm,xshift=0.5cm]C)--([yshift=-0.5cm,xshift=0.5cm]C)-- ([yshift=-0.5cm,xshift=-0.5cm]D)--([yshift=0.5cm,xshift=-0.5cm]D)--([yshift=0.5cm,xshift=0.5cm]D)--([yshift=-0.5cm,xshift=0.5cm]A)--([yshift=-0.5cm,xshift=-0.5cm]A)--([yshift=0.5cm,xshift=-0.5cm]A);
		
		
		\draw ([yshift=-0.3cm,xshift=0.3cm]A)-- ([yshift=-0.3cm,xshift=-0.3cm]B)--
		([yshift=0.3cm,xshift=-0.3cm]B) --([yshift=0.3cm,xshift=0.3cm]B)--([yshift=-0.3cm,xshift=0.3cm]C)--([yshift=-0.3cm,xshift=-0.3cm]C)--([yshift=0.3cm,xshift=-0.3cm]C)-- ([yshift=0.3cm,xshift=0.3cm]D)--([yshift=-0.3cm,xshift=0.3cm]D)--([yshift=-0.3cm,xshift=-0.3cm]D)--([yshift=0.3cm,xshift=-0.3cm]A)--([yshift=0.3cm,xshift=0.3cm]A)--([yshift=-0.3cm,xshift=0.3cm]A);

	\end{tikzpicture}
	\begin{center}
		\begin{tabular}{l*{40}{@{\hskip.05mm}c@{\hskip.8mm}} c c}
			\begin{tabular}{c}
				
		\footnotesize{\textbf{R\'EPUBLIQUE DU CAMEROUN}} \\
				
				\scriptsize{\textbf{****************}} \\
				
					\scriptsize{\textbf{Paix - Travail - Patrie}} \\
				
			\scriptsize{\textbf{******************}}\\ 
			\footnotesize{	\textbf{UNIVERSIT\'E DE YAOUND\'E I}}\\
				
			\scriptsize{	\textbf{****************}} \\
				
			\footnotesize{	\textbf{ECOLE NATIONALE SUPERIEURE}} \\
			\footnotesize{	\textbf{POLYTECHNIQUE DE YAOUNDE}} \\
				
			\scriptsize{	\textbf{****************}} \\
		   \scriptsize{	\textbf{D\'EPARTEMENT DE GENIE}}\\
		   \scriptsize{	\textbf{INFORMATIQUE}}\\
				
			\scriptsize{	\textbf{****************}}\\
				
			\end{tabular} &
			\begin{tabular}{c}
				
				\includegraphics[height=4cm, width=2.8cm]{po.png}
				
			\end{tabular} &
			\begin{tabular}{c}
				
				\footnotesize{\textbf{ REPUBLIC OF CAMEROON}} \\
				
				\footnotesize{\textbf{****************}} \\
				
					\scriptsize{\textbf{Peace - Work - Fatherland}} \\
				
				\scriptsize{\textbf{****************}} \\
				\footnotesize{\textbf{UNIVERSITY OF YAOUNDE I}}\\
				
				\scriptsize{\textbf{****************}} \\
				
				\footnotesize{\textbf{NATIONAL ADVANCED SCHOOL}} \\
				\footnotesize{\textbf{OF ENGINEERING OF YAOUNDE}} \\
				
				\scriptsize{\textbf{****************}} \\
				\scriptsize{\textbf{DEPARTMENT OF COMPUTER}}\\
				\scriptsize{\textbf{ENGINEERING}}\\
				
				\footnotesize{\textbf{****************}}\\
				
			\end{tabular}	
		\end{tabular}
	
		\vspace{0.5cm}
		\begin{tabular}{l*{40}{@{\hskip 3.5cm}c@{\hskip5cm}} p{3.5cm} r}
		\end{tabular}
		
		\noindent\rule{\textwidth}{0.7mm}
		\Large{{\textbf{CHAPITRE 1}}}\\
		\Large{{\textbf{\textit{EXERCICES}}}}
		\noindent\rule{\textwidth}{0.7mm}
	\end{center}
		
	\begin{center}
	\begin{tabular}{c}
		
		\vspace{0.1cm}
		\normalsize
	
	
		\vspace{0.1cm}
		\normalsize\textbf{Option }:\\			
		\textsl{Cybersécurité et Investigation Numérique}
		
	\end{tabular}
	\end{center}
		
	\begin{center}
		\normalsize %\hspace{-2cm}
		\begin{tabular}{c}
			\vspace{0.07cm}
			\hspace{0.02cm} \textbf{\textbf{Rédigé par :}}\\
			\hspace{0.02cm} \textsl{\textbf{NGWAMBE  MARIELLA}, 22P040}\\\\
			
		\end{tabular}
	\end{center}
	
	\begin{center}
	\hspace{0.02cm} \textbf{Sous l'encadrement de:}\\
	\hspace{0.02cm} \textsl{M. Thierry MINKA}
	\end{center}
	
    
	\vspace{2cm}
	\begin{center}
		\textbf{Année académique 2025 / 2026}
	\end{center}
		
	\vspace{-1.4cm}
	
		
	\vfill%\null
	
\end{titlepage}
\tableofcontents
\newpage
\begin{document}
 \chapter{} PARTIE 1: 
   \section{ Exercice 1}
 \paragraph{} — Paradoxe de la transparence}
Dans notre société contemporaine, la transparence est souvent présentée comme une valeur incontournable, synonyme d’authenticité, de vérité et de confiance. Pourtant, le philosophe Byung-Chul Han propose une lecture plus critique de ce concept. Dans son essai La société de la transparence, il affirme que la transparence « neutralise tout dans une distance uniforme ». Cette citation, à la fois dense et provocatrice, interroge les effets profonds de la transparence sur nos relations,nos échanges et même notre perception du réel. Elle invite à repenser ce que nous valorisons comme "clarté", en soulignant que l'exigence de transparence pourrait paradoxalement appauvrir la richesse du lien humain et de la complexité sociale. La « société de la transparence » décrite par Byung-Chul Han met en lumière
une tension essentielle : la volonté de tout rendre visible pour protéger la
démocratie et lutter contre la corruption, contre le droit fondamental à l'intimité.
Dans l'espace numérique, onde de traçabilité et d'indexation, la transparence
fonctionne comme remède et poison : elle améliore la reddition de comptes
mais fragilise la sphère privée, creuse les asymétries d'information et
expose les individus à des formes nouvelles de violence symbolique.
Considérée sous l'angle de l'investigation numérique, cette tension se traduit
par un conflit entre deux impératifs légitimes : la recherche de preuves et la
protection des données personnelles.

Un cas paradigmatique est l'accès aux logs d'un service public dans le cadre
d'une enquête sur la corruption. L'ouverture de ces traces permet de
reconsitituer les chaînes d'événements, d'identifier des acteurs et de
garantir la transparence administrative. En revanche, la divulgation de ces
logs risque d'exposer des citoyens innocents, de révéler des habitudes
sensible ou des éléments protégés par le secret professionnel. Le risque n'est
pas seulement individuel : l'obsession de la traçabilité érode la confiance,
puisque les usagers modifient leurs comportements pour éviter d'être
surveillés, limitant l'accès à l'information et appauvrissant l'espace public.

L'éthique kantienne offre un cadrage normatif pertinent. Le principe moral
fondamental kantien — agir de telle sorte que l'humanité soit toujours
traitée comme une fin et jamais seulement comme un moyen — impose
une contrainte sur l'usage de la transparence. Appliqué à l'investigation
numérique, ce principe exige que la collecte, la conservation et la divulgation
des preuves ne traitent jamais la personne comme simple ressource
informationnelle. Concrètement, cela se traduit par des règles de proportionnalité,
de minimisation des données, d'anonymisation lorsque possible, et d'un
contrôle judiciaire effectif pour autoriser l'accès aux informations sensibles.

Une résolution pratique combine plusieurs leviers : (i) \emph{principe de
minimisation} : ne collecter que ce qui est strictement nécessaire ; (ii)
\emph{garde-fou procédural} : autorisation judiciaire préalable pour les
accès intrusifs ; (iii) \emph{techniques de protection} : chiffrement, zoning,
masquage et protocoles zero-knowledge pour prouver sans révéler ; (iv)
\emph{transparence procédurale} : audit des accès et droit d'information
pour les personnes affectées. Ainsi la recherche de la vérité reste possible
sans sacrifier la dignité humaine.

En conclusion, la solution n'est ni l'absolue transparence ni le secret total,
mais un équilibre institutionnalisé : une transparence encadrée, respectueuse
des droits, qui fait de l'investigateur non un voyeur mais un gardien de la
déontologie. Ce compromis incarne la lecture kantienne : protéger la société
tout en respectant l'individu comme une fin en soi.
  \section*{Exercice 2 — Transformation ontologique}
    \subsection*{Comparaison Heidegger / ère numérique}
\paragraph Heidegger conçoit la technique non simplement comme un outil, mais comme un mode de révélation du monde – ce qu’il nomme Gestell (l’arraisonnement). La technique moderne, selon lui, transforme notre rapport à l’être en réduisant le réel à un ensemble de ressources disponibles, exploitables et calculables. Ce processus modifie profondément le Dasein, l’être humain compris comme être-au-monde, c’est-à-dire en relation ouverte, vécue et située avec le monde qui l’entoure. À l’ère numérique, cette pensée trouve une résonance nouvelle : la technique ne se contente plus d’outiller ou d’optimiser l’action humaine, elle reconfigure l’existence même. Le Dasein contemporain ne vit plus seulement dans un monde physique, mais il est augmenté ou dédoublé par une présence numérique continue : profils, historiques, traces, logs… autant d’extensions du soi qui persistent indépendamment de la conscience immédiate, et qui façonnent une nouvelle forme d’« être-là » dans le monde digital.Mais ce double numérique, loin d’enrichir nécessairement l’expérience, peut participer à un appauvrissement existentiel, comme le craignait Heidegger. Le monde numérique tend à réduire l’existence à des données, des comportements mesurables, des préférences exploitables – un monde indexé, quantifié et monétisé. Ce qui était vécu dans l’opacité, la temporalité et l’ouverture devient ainsi transparent, archivé et contrôlable, réactualisant la critique heideggérienne : une révélation qui dissimule, une lumière technique qui occulte la profondeur de l’être.

\subsection*{Analyse d'un profil social comme \emph{être-par-la-trace}}
Un profil social constitue une manifestation intentionalissee de l'individu. Il ne s'agit pas seulemeent d'un ensemble de donnees techniques, mais d'un espace d'expression de l'existence ou se materialisent les choix, les preferences , les relations et les habitudes. l'analyse hermeneutiaue dun profil suppose de le lire comme un texte : chaque post, commentaire ou photo ne prend sens au' en assemblant tous ces elements. l'etre-par-la-trace traduit ainsi la tension entre la presence et la representation.C’est dans cette perspective que l’on identifie cinq types de traces existentielles révélant la présence numérique de l’être : les relations sociales (amis, abonnés, followers) traduisent les liens intersubjectifs et la manière dont l’individu se situe dans un réseau de relations ; les centres d’intérêt (pages suivies, contenus likés ou commentés) expriment les préférences, valeurs et affinités symboliques ; la temporalité (fréquence de connexion, rythme de publication) manifeste un rapport singulier au temps numérique, où la continuité, l’absence ou la ponctualité deviennent signifiantes ; la spatialité (géolocalisations, lieux mentionnés) inscrit l’existence dans des espaces réels et virtuels, ancrant le corps dans le monde ; enfin, l’intentionalité (contenus partagés) révèle la direction du sens donnée à l’action : se montrer, témoigner, influencer ou simplement exister. Ces traces, en se combinant, donnent corps à un être-par-la-trace, où l’identité se manifeste et se légitime sans jamais se réduire à la donnée technique. Cette approche transforme également la notion de preuve légale : la preuve n’est plus un simple objet statique, mais un processus relationnel reliant la trace à l’identité et à l’intention de son auteur. La valeur probante naît alors de la cohérence herméneutique entre le signe et le sens, entre le sujet et ses traces. La validation ontologique de cette lecture repose sur la reconnaissance de la dimension d’être dissimulée derrière la technique : chaque donnée, chaque interaction numérique, devient le signe d’un être-en-relation, d’un être-temps et d’un être-espace. Le profil social s’impose dès lors comme une véritable cartographie existentielle, où l’identité se joue entre le visible (la trace) et l’invisible (le sens).

\subsection*{Impact sur la preuve légale}
Les preuves numériques sont immatérielles, volatiles et manipulables.
La transformation ontologique introduit trois défis :
\begin{itemize}
  \item \textbf{Stabilité} : copies binaires et versions multiples ;
  \item \textbf{Authenticité} : garantir origine et intégrité (hash, chain of custody) ;
  \item \textbf{Interprétation} : nécessité d'une herméneutique technique pour éviter
  des conclusions hâtives (contexte, corrélations vs causalité).
\end{itemize} 
  PARTIE 2: 
   \section{Exercice 3} 
   \paragraph{}Nous allons calculer l'entropie de chaque fichier 
\begin{lstlisting}
import math
from collections import Counter
def entropy(filepath):
    with open(filepath, 'rb') as f:
        data = f.read()
    if not data:
        return 0.0
    counts = Counter(data)
    n = len(data)
    H = 0.0
    for c in counts.values():
        p = c / n
        H -= p * math.log2(p)
    return H  # bits par octet

# Usage
 print(entropy('6hearty.txt'))
 print(entropy('OIP (2).jpeg'))
 print(entropy('document.pdf.aes'))
\end{lstlisting}
Voici les résultats obtenus respectivement: 
3.9362600275315254,7.96551329191923,5.9999426266173135
D'après les attentes : on peut conclure que: 
\begin{item}
    \item -Le fichier texte est conforme car son entropie est proche de la borne haute 
    \item -Le fichier image légèrement hors de la borne ne présente rien d'anormal car certains JPEG très compressés comme dans notre cas peuvent atteindre cette valeur 
    \item -le résultat obtenu (très faible) suggère que peut contenir beaucoup d'en-tête non chifffrées par exemple.
\end{item}
Le seuil proposé H<7.7. Si ce n'est pas le cas on peut déduire qu'il s'agit d'un chiffrement suspect ou données chiffrées.
\section*{Exercice 5:Modélisation de l'Effet Papillon en Forensique }\paragraph{Modélisation mathématique} On part de la relation exponentielle :

[
\delta(t) = \delta(0)e^{\lambda t}
\]

avec \(
\delta(0) = 30~s
\)

δ(0)=30 s (modification initiale). Pour estimer \delta(t)
λ \begin{enumerate} \item On mesure les décalages absolus observés $\delta(t_i)$ pour les 
��
i `events` après la perturbation. \item On construit le nuage $(t_i,\ln\delta(t_i))$ et on effectue une régression linéaire : 
ln
⁡
��
(
��
)
≈
��
 
��
+
��
lnδ(t)≈at+b alors 
��
^
=
��
λ
^
=a et 
ln
⁡
��
(
0
)
≈
��
lnδ(0)≈b. \end{enumerate}

\paragraph{Exemple numérique (simulation réalisée)} Avec une simulation calibrée on peut choisir 
��
��
��
��
��
=
0.005
 
s
−
1
λ
true
	​

=0.005s
−1
. La régression sur les données bruitées fournit une estimée 
��
��
��
��
≈
0.00500035
 
s
−
1
λ
est
	​

≈0.00500035s
−1
 avec un excellent R
2
2
 (≈0.9997). Cela illustre que : \begin{itemize} \item même une petite perturbation (30 s) peut, si le système est instable (
��
>
0
λ>0), produire des décalages significatifs \item l'estimation de 
��
λ permet de définir un horizon temporel de confiance : pour un seuil acceptable $\Delta_{max}$, on résout pour $t$ tel que $\delta(0)e^{\lambda t}=\Delta_{max}$, d'où $t=\frac{1}{\lambda}\ln(\Delta_{max}/\delta(0))$. \end{itemize}

\section{Partie 3 — Révolution quantique et implications}

\subsection{Exercice 6 : Chat de Schrödinger adapté}

Un fichier avant analyse peut être modélisé comme un 'quasi-superposé' entre intact'' et effacé'' : 
∣
��
⟩
=
��
∣
0
⟩
+
��
∣
1
⟩
∣ψ⟩=α∣0⟩+β∣1⟩ avec $|\alpha|^2$ la probabilité que la preuve soit valide. En pratique, la mesure correspond à l'opération d'extraction de métadonnées et peut altérer l'objet (ex. mesures destructives). Un protocole pragmatique consiste \begin{enumerate} \item Échantillonnage non-destructif (hashing, métadonnées) ; \item mesures faibles si disponible (techniques de forensique non invasive) ; \item redondance multi-source pour confirmer sans mesurer chaque artefact. \end{enumerate}

\subsection{Exercice 7 : Calculs sur la sphère de Bloch}

Pour $\theta=\pi/3$ et $\phi=\pi/4$ :

\[
|\psi\rangle 
= \cos\!\left(\frac{\theta}{2}\right)|0\rangle 
+ e^{i\phi}\sin\!\left(\frac{\theta}{2}\right)|1\rangle
= \frac{\sqrt{3}}{2}|0\rangle 
+ \frac{1}{2}e^{i\pi/4}|1\rangle
\]

Les probabilités de mesure sont :
\[
P(0) = \left(\frac{\sqrt{3}}{2}\right)^2 = 0.75, 
\quad
P(1) = \left(\frac{1}{2}\right)^2 = 0.25
\]

Coordonnées sur la sphère de Bloch :
\[
x = \sin\theta \cos\phi = \sin\!\left(\frac{\pi}{3}\right)\cos\!\left(\frac{\pi}{4}\right) \approx 0.612,
\quad
y = \sin\theta \sin\phi \approx 0.612,
\quad
z = \cos\theta = 0.5
\]

y=sinθsinϕ≈0.612,z=cosθ=0.5


Implication : avant mesure, la probabilité d'obtenir la preuve valide est 75 pourcents, ce qui impose de raisonner en termes de confiances et non de certitudes absolues.

\subsection{Exercice 8 : Théorème de non-clonage}

La démonstration classique (par contradiction) montre l'impossibilité d'un opérateur linéaire clonant les 'qubits'. L'implication forensique est que l'on ne peut pas réaliser des copies fidèles d' états quantiques inconnus : la conservation des preuves exige des protocoles alternatifs (engagements, preuves ZK, mesures non destructives codifiées) et une nouvelle chaîne de custody quantique.

\section*{Partie 4 — Paradoxe de l'Authenticité Invisible}

\subsection{Exercice 9 : Formalisation mathématique (exemple)}

Choisissons trois systèmes et leur notation approximative sur $[0,1]$ : \begin{itemize} \item Signature RSA classique : $A=0.9,; C=0.2,; O=0.85$. \item ZK-SNARK : $A=0.75,; C=0.85,; O=0.4$. \item ZK-STARK : $A=0.8,; C=0.8,; O=0.6$. \end{itemize}
\noindent Pour rendre homogène la notation avec l'énoncé qui demande $A\cdot C \le 1-\delta$, nous posons la correspondance :
\[
A \triangleq R \quad\text{(on identifie ici la composante $R$ fournie par la méthodologie à l'authenticité }A).
\]

\subsection*{Calculs élémentaires}
Calculons $A\cdot C$ pour chaque système :
\[
\begin{aligned}
\text{RSA: } & A_{RSA}=0.9,\quad C_{RSA}=0.1,\quad A_{RSA}\cdot C_{RSA}=0.9\times 0.1=0.09,\\[4pt]
\text{ZK-SNARK: } & A_{SNARK}=0.7,\quad C_{SNARK}=0.9,\quad A_{SNARK}\cdot C_{SNARK}=0.7\times 0.9=0.63,\\[4pt]
\text{ZK-STARK: } & A_{STARK}=0.8,\quad C_{STARK}=0.8,\quad A_{STARK}\cdot C_{STARK}=0.8\times 0.8=0.64.
\end{aligned}
\]

La borne $\delta$ minimale qui assure l'inégalité $A\cdot C \le 1-\delta$ pour \emph{tous} les systèmes est :
\[
\delta_{\min} = 1 - \max_{\text{sys}}(A\cdot C) = 1 - 0.64 = 0.36.
\]
Autrement dit, pour garantir l'inégalité sur les trois systèmes, il faut $\delta \le 0.36$.

\subsection*{Estimation de $\hbar_{\text{num}}$}
Nous devons estimer empiriquement une constante analogue $\hbar_{\text{num}}$ à partir de l'ensemble des systèmes, en utilisant une forme heuristique de l'inégalité d'incertitude :
\[
\Delta A \cdot \Delta C \ge \frac{\hbar_{\text{num}}}{2}.
\]
Pour produire une estimation pratique, nous adoptons l'hypothèse simple (heuristique) suivante pour les \emph{incertitudes} :
\[
\Delta A \approx 1 - A,\qquad \Delta C \approx 1 - C,
\]
c'est-à-dire que l'incertitude relative sur chaque axe est mesurée par l'écart à la valeur maximale \(1\). Calculons $\Delta A\cdot \Delta C$ pour chaque système :
\[
\begin{aligned}
\text{RSA: } & \Delta A = 1-0.9=0.1,\quad \Delta C=1-0.1=0.9,\quad \Delta A\Delta C=0.09,\\[4pt]
\text{ZK-SNARK: } & \Delta A = 1-0.7=0.3,\quad \Delta C=1-0.9=0.1,\quad \Delta A\Delta C=0.03,\\[4pt]
\text{ZK-STARK: } & \Delta A = 1-0.8=0.2,\quad \Delta C=1-0.8=0.2,\quad \Delta A\Delta C=0.04.
\end{aligned}
\]

Moyennant l'hypothèse que l'inégalité moyenne tient au bord (approximation), nous estimons :
\[
\overline{\Delta A\Delta C} = \frac{0.09+0.03+0.04}{3} \approx 0.05333.
\]
D'après $\Delta A\Delta C \gtrsim \hbar_{\text{num}}/2$ nous obtenons comme estimation prudente :
\[
\boxed{\hbar_{\text{num}} \approx 2\times \overline{\Delta A\Delta C} \approx 0.11.}
\]

\paragraph{Remarques sur l'interprétation}
\begin{itemize}
  \item Cette estimation est \textbf{heuristique} : la valeur numérique dépend fortement du choix de la métrique d'incertitude ($\Delta A,\Delta C$) et du petit échantillon (3 systèmes). D'autres définitions de $\Delta$ (écart-type, intervalle de confiance) donneraient des résultats différents.
  \item Le guide mentionnait des résultats typiques $\hbar_{\text{num}}\approx 0.3$--$0.5$. la différence s'explique par la métrique d'incertitude : si l'on adopte une mesure d'incertitude plus conservatrice (p.ex. $\Delta A = 0.5(1-A)$ ou en utilisant erreurs empiriques issues d'une calibration), l'estimation peut augmenter et rejoindre l'intervalle 0.3--0.5.
  \item Conclusion opérationnelle : il est préférable de calibrer $\hbar_{\text{num}}$ sur un large jeu de systèmes et d'adopter une définition standardisée de $\Delta$ (p.ex. écart-type des évaluations d'experts) avant de tirer des conclusions normatives.

\vspace{6pt}
\noindent\rule{\textwidth}{0.4pt}

\section*{Exercice 10 — Implémentation simplifiée ZK-NR (proof-of-concept)}
\subsection*{Objectif}
Fournir un prototype pédagogique en Python qui simule les étapes architecturales : engagement Merkle + preuve (ici simulée) + vérification. Mesurer empiriquement le compromis confidentialité / vérifiabilité et l'overhead.

\subsection*{Code (proof-of-concept simplifié)}
\begin{verbatim}
# simple_zknr.py -- proof-of-concept pédagogique (non sécurisé cryptographiquement)
import hashlib
import time
import math

class MerkleTree:
    def __init__(self, leaves):
        self.leaves = [self._leaf_hash(l) for l in leaves]
        self.root = self._build_root(self.leaves)

    def _leaf_hash(self, data):
        return hashlib.sha256(data).hexdigest()

    def _build_root(self, leaves):
        nodes = leaves[:]
        if not nodes:
            return None
        while len(nodes) > 1:
            if len(nodes) % 2 == 1:
                nodes.append(nodes[-1])
            pairs = []
            for i in range(0, len(nodes), 2):
                pairs.append(hashlib.sha256((nodes[i]+nodes[i+1]).encode()).hexdigest())
            nodes = pairs
        return nodes[0]

    def get_root(self):
        return self.root

class DummyZKProver:
    # Cette classe simule une "preuve" : en réalité ce n'est pas une ZK-proof.
    def prove(self, statement, witness, commitment):
        # Simuler un temps de génération O(n log n) approximé par n*log(n)
        n = len(witness)
        ops = n * math.log(max(2, n))
        time.sleep(min(0.01, ops*1e-6))  # temps simulé très court
        # création d'une "preuve" factice (hash)
        s = statement.encode() + b'|' + commitment.encode()
        return hashlib.sha256(s).hexdigest()

    def verify(self, proof, statement, commitment):
        s = statement.encode() + b'|' + commitment.encode()
        expected = hashlib.sha256(s).hexdigest()
        return proof == expected

class SimpleZKNR:
    def __init__(self):
        self.merkle_tree = None
        self.prover = DummyZKProver()

    def create_proof(self, data_list, statement):
        # Etape 1 : engagement Merkle
        mt = MerkleTree(data_list)
        commitment = mt.get_root()
        # Etape 2 : preuve (simulée)
        proof = self.prover.prove(statement, data_list, commitment)
        return commitment, proof

    def verify_proof(self, commitment, proof, statement):
        return self.prover.verify(proof, statement, commitment)

# --- Mesure d'overhead expérimental (exemple) ---
if __name__ == "__main__":
    for n in [100, 1000, 5000]:
        data = [("item"+str(i)).encode() for i in range(n)]
        z = SimpleZKNR()
        start = time.perf_counter()
        commitment, proof = z.create_proof(data, "statement")
        gen_t = time.perf_counter() - start
        start = time.perf_counter()
        ok = z.verify_proof(commitment, proof, "statement")
        ver_t = time.perf_counter() - start
        print(f"n={n}: gen_time={gen_t:.4f}s, ver_time={ver_t:.6f}s, proof_ok={ok}")
\end{verbatim}

\subsection*{Explication et limites}
\begin{itemize}
  \item Le \texttt{MerkleTree} ci-dessus calcule un engagement de racine Merkle : coût de construction $\mathcal{O}(n)$ en pratique (avec concaténations et hachages), taille d'engagement constante (la racine).
  \item \texttt{DummyZKProver} ne fournit pas de vraie preuve ZK : il sert à simuler des coûts et des métriques. Pour un prototype réel, on intègre une librairie (libsnark, arkworks, circom + snarkjs, ou une implémentation STARK).
  \item La fonction \texttt{prove} simule un coût croissant (heuristique $n\log n$) et retourne une preuve représentée par un hash ; la vérification recompute un hash attendu (temps $\approx O(1)$).
\end{itemize}

\subsection*{Mesures attendues et interprétation}
\begin{itemize}
  \item \textbf{Temps de génération} : prototype vrai ZK-SNARK/STARK $\sim O(n\log n)$ (dépend du circuit, du pre-processing). Dans le PoC ci-dessus la génération croît avec $n$ (mesures imprimées par le script).
  \item \textbf{Temps de vérification} : pour SNARKs/STARKs bien conçus la vérification est très rapide (souvent quasi-constante ou polylog), d'où l'approximation $O(1)$ utile.
  \item \textbf{Taille de la preuve} : SNARK $\approx O(1)$ à $O(\log n)$ suivant systèmes ; STARKs tendent vers $O(\log n)$ ou plus, mais sans trusted setup.
  \item \textbf{Overhead acceptable} : l'exigence donnée (<30\% du temps total) se mesure ainsi :
    \[
    \text{overhead} = \frac{t_{\text{ZK}}}{t_{\text{opération totale}}}.
    \]
    Pour un workflow d'enquête, mesurer $t_{\text{ZK}}$ (génération + vérification) vs temps total de traitement permet d'accepter ou non le recours ZK.
\end{itemize}

\subsection*{Protocole d'expérience conseillé}
\begin{enumerate}
  \item Générer jeux de données synthétiques pour $n\in\{100,1000,5000,10000\}$.
  \item Exécuter le PoC et enregistrer $t_{\text{gen}}(n)$, $t_{\text{ver}}(n)$ et taille(proof).
  \item Calculer overhead relatif pour un scénario réaliste (p.ex. vérification documentaire complète) et vérifier la contrainte «~$<$ 30\%~».
  \item Si overhead trop élevé, explorer : réduction de $n$ via agrégation, utilisation de proofs récursives, ou déporter la génération hors-chaîne.
\end{enumerate}

\subsection*{Conclusion (exercice 10)}
Le proof-of-concept fourni est pédagogique et permet :
\begin{itemize}
  \item de mesurer empiriquement des métriques de coût,
  \item d'étudier le compromis confidentialité / vérifiabilité (en jouant sur la granularité des engagements et la quantité d'information incluse dans le témoin),
  \item d'estimer si l'overhead computationnel reste dans la tolérance pratique (<30\%).
\end{itemize}
 \section*{Partie 5: Intégration et Synthèse Avancée}
 \subsection{Exercice 11 : Étude de cas \emph{QuantumLeaks}}

\paragraph{Contexte}
L’affaire \emph{QuantumLeaks} illustre les nouveaux défis posés à la justice numérique : fuites de données chiffrées, incertitudes liées à la temporalité des journaux, et nécessité d’authentifier sans violer la confidentialité. L’objectif est de concevoir une architecture d’investigation résiliente, respectueuse des principes éthiques et compatible avec la cryptographie post-quantique.

\paragraph{Analyse du cas}
Trois axes principaux structurent la réflexion :
\begin{enumerate}
  \item \textbf{Authenticité} : garantir que les preuves numériques ne soient pas altérées — recours à des signatures post-quantiques telles que \textit{CRYSTALS-Dilithium}.
  \item \textbf{Confidentialité} : préserver les données sensibles à travers des preuves à divulgation nulle de connaissance (ZK-STARKs) permettant de vérifier sans exposer.
  \item \textbf{Traçabilité} : assurer une chaîne de custody transparente via des registres immuables (blockchain ou ledgers audités).
\end{enumerate}

\paragraph{Proposition de solution intégrée}
\begin{itemize}
  \item \textbf{Architecture hybride} : combiner les signatures Dilithium (authenticité) avec des preuves ZK-STARK (confidentialité), articulées dans un registre auditable (ledger Merkle horodaté).
  \item \textbf{Stockage sécurisé} : recourir à un chiffrement hétérogène (AES + post-quantique) avec rotation périodique des clés et archivage hors ligne certifié.
  \item \textbf{Chaîne de custody} : journalisation immuable, réplication multi-niveaux, preuves d’intégrité signées par des tiers de confiance.
  \item \textbf{Conformité juridique} : application du RGPD, encadrement légal des horodatages et de l’anonymisation, validation procédurale des preuves.
\end{itemize}

\paragraph{Conclusion}
Le cas \emph{QuantumLeaks} met en évidence le basculement d’une approche statique de la preuve à une approche \textbf{processuelle et systémique}, où chaque couche (cryptographique, juridique, ontologique) co-construit la validité. La philosophie sous-jacente rejoint celle de la «~preuve relationnelle~» : une vérité numérique se fonde moins sur la possession que sur la cohérence entre ses relations.

\subsection{Exercice 12 : Débat philosophique — Réalisme vs Constructivisme}

\paragraph{Problématique}
La question de la neutralité de l’observateur en investigation numérique oppose deux perspectives :  
\textbf{le réalisme scientifique}, pour qui les artefacts existent indépendamment de l’observateur, et  
\textbf{le constructivisme}, pour qui la connaissance est co-produite par le dispositif technique et l’analyste.

\paragraph{Position réaliste}
Les réalistes affirment que les traces numériques sont des faits bruts, objectivement mesurables. Dans cette perspective, les outils de mesure (hash, logs, métadonnées) ne font que révéler ce qui \emph{est déjà là}. L’observateur agit comme un miroir : neutre, rationnel, garant de la vérité factuelle. Référence : Wheeler et la notion d’«~it from bit~», où l’information précède l’interprétation.

\paragraph{Position constructiviste}
Les constructivistes, inspirés de Kuhn et Latour, rappellent que toute observation est médiée par un paradigme technique. Les outils (logiciels, algorithmes, filtres d’analyse) configurent ce qui est perçu. La preuve numérique n’est donc pas découverte, mais construite. L’interprétation, la sélection des indicateurs et la temporalité d’analyse influencent la vérité produite.

\paragraph{Synthèse critique}
La position de synthèse s’appuie sur le modèle \textbf{CRO (Confidentialité–Rationalité–Ouverture)} :  
chaque acte d’investigation doit équilibrer la rigueur de la mesure (réalisme) et la réflexivité de l’interprète (constructivisme). La neutralité absolue est un idéal régulateur, non une réalité empirique.

\paragraph{Conclusion}
L’analyste forensique contemporain n’est ni pur spectateur ni simple technicien, mais un \textit{médiateur de vérité numérique}. Son rôle éthique consiste à maintenir la cohérence entre technique, droit et sens, dans le respect du principe ontologique d’«~être-par-la-trace~».

\subsection{Exercice 13 : Projet de recherche personnel}

\paragraph{Titre proposé}
\emph{Mesure de l’impact de la superposition quantique sur l’admissibilité des preuves numériques.}

\paragraph{Problématique}
Les états quantiques ne peuvent être copiés fidèlement (théorème de non-clonage). Cette propriété remet en question la reproductibilité, pilier de la preuve scientifique et judiciaire. Comment définir des critères d’admissibilité pour des preuves issues de systèmes quantiques instables ou superposés ?

\paragraph{Hypothèse}
La superposition et l’effet de mesure imposent un seuil de confiance minimal, au-delà duquel la preuve perd son statut d’objectivité. Ce seuil dépend de la topologie du système et du protocole de mesure adopté.

\paragraph{Méthodologie proposée}
\begin{enumerate}
  \item \textbf{Simulation} : génération de données quantiques simulées (via Qiskit ou QuTiP) intégrant des mesures faibles et des états mixtes.
  \item \textbf{Modélisation statistique} : application d’un modèle de type Monte Carlo pour estimer la variance des résultats de mesure et définir un intervalle de confiance pour chaque état.
  \item \textbf{Évaluation légale} : comparaison des seuils obtenus aux exigences juridiques d’admissibilité de la preuve (intégrité, traçabilité, reproductibilité).
\end{enumerate}

\paragraph{Livrables attendus}
\begin{itemize}
  \item Rapport scientifique et éthique sur la «~preuve quantique~».
  \item Scripts reproductibles en Python (Qiskit/NumPy).
  \item Visualisations de la dispersion des mesures et zones d’incertitude.
  \item Recommandations pour la normalisation légale de la preuve quantique.
\end{itemize}

\paragraph{Conclusion générale de la Partie 5}
Cette dernière partie opère la jonction entre technique, droit et ontologie. L’investigation numérique devient une pratique \textbf{interdisciplinaire et réflexive} : elle ne cherche plus seulement à prouver, mais à comprendre les conditions de possibilité de la preuve elle-même.  
Ainsi se clôt la boucle : de la trace technique à la vérité relationnelle, l’investigation devient une \emph{phénoménologie appliquée du numérique}.
\end{document}
